\section{2d/3d 对偶：CS/WZW 与 AdS$_3$/CFT$_2$ 是两码事}
\label{sec:duality}

\subsection{这一章要解决什么}

本文的标题里有“量子场论”，正文里反复出现“三维理论”和“二维共形场论”，读者很容易顺手把它归入后来流行的“AdS$_3$/CFT$_2$ 全息对偶”。这个归类是错的，而且错在两个层次上。第一，本文的三维--二维关系是\emph{态空间的等同}：三维 Chern--Simons 理论在曲面 $\Sigma$ 上的 Hilbert 空间，就是二维共形块空间；原文用 “are” 这个词，说的是同一个向量空间的两种描述。第二，AdS$_3$/CFT$_2$ 是\emph{两个动力学理论之间的等价}：三维量子引力（含度规、有引力子、有黑洞）与边界二维 CFT 的配分函数相等，靠的是一整套边界条件与全息字典。本章要把这两件事分开，逐条对照它们的体自由度、边界条件、中心荷、可证性，并说明本文与“全息”唯一的接触点——“三维引力可以写成 Chern--Simons 理论”——为什么只是经典改写，不是等同。

\subsection{原文说了什么}

\paragraph{（一）本文的三维理论没有度规。}
原文在选作用量时明确说：\emph{“we want to pick a Lagrangian which does not require any choice of metric”} \src{353}。给出的作用量是
\begin{equation}
L=\frac{k}{4\pi}\int_M\Tr\Big(A\wedge \dd A+\frac23A\wedge A\wedge A\Big), \tag{1.3}
\end{equation}
它只需要定向，不需要度规。与 (1.2) 的 Yang--Mills 作用量对比：那里的 $\sqrt g\,g^{ik}g^{jl}$ 必须用度规缩并，而 (1.3) 的积分是 3-形式的积分。没有度规，就没有引力子，也没有“距离”“面积”“黑洞熵”这些概念。这一条是后面对照的出发点。

\paragraph{（二）态空间是有限维的。}
原文在 \S3 论证：约束之后，相空间是平坦联络模空间，它是有限维的。原文说这些空间\emph{“will turn out to be finite dimensional spaces”} \src{366}，又称之为\emph{“eminently finite dimensional phase space”} \src{367}。有限维是“拓扑场论”的标志：一个三维理论若有传播的局部自由度，其态空间一般不是有限维的。本文的 $\cH_\Sigma$ 有限维，说明它没有体自由度。

\paragraph{（三）二维--三维的等式是一条态空间陈述。}
原文 p.370 的英文原句是：
\begin{quote}
\emph{“the space of conformal blocks in 1 + 1 dimensions are the quantum Hilbert spaces obtained by quantizing a 2 + 1 dimensional theory.”}
\end{quote}
\src{370} 注意动词是 “are”，不是 “correspond to” 或 “are dual to”：原文主张的是\emph{同一个空间}，不是两个理论之间的映射。将来到 p.396 总结时，原文的措辞变成 “can be interpreted as”（可以解释为），语气明显弱了一档。两处措辞不同不是笔误：p.370 说的是数学对象的等同（同一个向量空间），p.396 说的是解释方式（把三维量子化得到的空间解释成共形块空间）。把 p.396 的 “can be interpreted as” 读成“体配分函数等于边界配分函数”，就把解释关系升级成了动力学等价。紧接着 p.371 说，Ward 恒等式把关联函数的确定归结为“在丛的模空间上选取 $\cL^{\otimes k}$ 的一个全纯截面”，而模空间上的射影平坦联络（有穿孔时是 Knizhnik--Zamolodchikov 联络）给出单值性。这一段的主语始终是“共形块空间”，不是“边界 CFT 的配分函数”。

\paragraph{（四）三维侧只在 $X\times S^1$ 上“接触”二维数据。}
本文给出的唯一一个把三维配分函数与二维 Hilbert 空间连起来的等式是
\begin{equation}
Z(X\times S^1)=\Tr_{\cH_X}(\mathbf 1)=\dim\cH_X . \tag{4.30}
\end{equation}
\src{387} 对 $X=S^2$ 且无标记点，(4.31) 给出 $Z(S^2\times S^1)=1$；有标记点时，(4.33)--(4.36) 把配分函数写成 $\cH_{S^2;\langle R\rangle}$ 的维数，特别是
\begin{equation}
Z(S^2\times S^1;R_a,R_b,R_c)=N_{abc}. \tag{4.36}
\end{equation}
\src{388} 这些等式全部是“三维配分函数 = 某个态空间的维数”，而不是“三维配分函数 = 二维配分函数”。一般三维流形（如 $S^3$）上的不变量是用手术从这些态空间拼出来的，仍然没有跳出“态空间”的范围。特别地取 $X=T^2$ 且无标记点：$\dim\cH_{T^2}=t$，$SU(2)$ 时 $t=k+1$，$SU(N)$ 时 $t=\binom{k+N-1}{N-1}$ \src{383--384}。这是有限维的，而且随 $k$ 只作多项式增长。这一点在 D-10.6 中会与“真正的二维 CFT Hilbert 空间”形成对照。

\paragraph{（五）原文对“联系”的自我定位。}
\S5 开头原文说：
\begin{quote}
\emph{“The basic connection that we have so far stated between general covariance in 2+1 dimensions and conformally invariant theories in 1+1 dimensions is that the physical Hilbert spaces obtained by quantization in 2+1 dimensions can be interpreted as the spaces of conformal blocks in 1+1 dimensions.”}
\end{quote}
\src{396} 这句话把“联系”限定得很死：\emph{Hilbert 空间可以被解释为共形块空间}。同页末尾说，这类考虑“促成了对 2+1 维引力的研究，其结果将在别处发表 [38]” \src{396}；p.397 又说“到这一步显然可以通过考虑更复杂的黎曼面与缝合……”（\emph{“It is obvious at this point that by considering more complicated Riemann surfaces with\dots”}）\src{397}。这两句是\emph{动机链}，不是等价性陈述。

\paragraph{（六）原文没有“全息”。}
对全译本做检索，“AdS”“holograph”“全息”“反德西特”没有任何实质命中（只有子串误匹配）。这不是疏忽：1988 年全息原理尚未提出，AdS/CFT  correspondence 要到 1997 年 Maldacena 的论文才出现。用 1997 年之后的框架去读 1988 年的文章，必须把“当时没有的东西”与“当时有的东西”分开。

\paragraph{（七）术语的对齐。}
原文用 “conformal blocks”（共形块）指 Ward 恒等式的解空间，用 “quantum Hilbert spaces” 指量子化三维理论得到的空间；p.370 的 “are” 把两者视为同一个对象。这里没有出现 “boundary CFT”“boundary condition”“holographic” 这些后来的术语。读本文时要把 “conformal block” 与 “boundary CFT partition function” 分开：前者是一个向量空间（对每个曲面 $\Sigma$），后者是一个数（对每个边界几何与边界耦合）。本章的所有对照都建立在这个区分上。

\subsection{原文没说的细节}

\begin{enumerate}[label=G10.\arabic*,leftmargin=*]
  \item \textbf{“对偶”这个词没有出现。}原文只说共形块空间“是”量子化三维理论得到的 Hilbert 空间，没有说两个理论“对偶”。本章沿用“CS/WZW 对应”这个说法，但把它严格定义为“态空间等同”，不把它当作动力学等价。\ours
  \item \textbf{本文没有边界条件的概念。}在闭流形上，CS 的态空间由 $\Sigma$ 的拓扑（亏格、标记点）决定，没有“渐近边界条件”这一层。AdS$_3$/CFT$_2$ 的核心恰是边界条件：固定渐近度规（或等价地固定两个 Virasoro 的生成元），才有无穷维渐近对称代数与 $c=3l/2G$。本文的有限维 $\cH_\Sigma$ 与“渐近对称的无穷维表示”是两种完全不同的结构。\ours
  \item \textbf{唯一接触点只占一句话。}“三维引力可以写成 Chern--Simons 理论”在本文只出现在 \S5 末尾的动机句与文献 [38] 中，没有证明，也没有用它推出任何关于二维理论的结论。原文没有说这个改写是“等同”，更没有说本文的紧群 CS 就是量子引力。\ours
  \item \textbf{1988 年的语境。}本文写作时既没有全息原理，也没有 AdS/CFT 字典。Brown--Henneaux 的 $c=3l/2G$ 虽在 1986 年已经出现，但它属于三维渐近对称分析，不是二维对偶；把它当作“本文的边界 CFT 中心荷”是时代错置（时间线见表 \ref{tab:timeline}）。本文能做的是“态空间”层面的严格对应：三维拓扑场论的模函子与二维 RCFT 的模函子相同（Segal 的框架，p.370--371），所以两边给出同一套 Hilbert 空间。这条对应比全息对偶更严格，但内容也更窄。\ours
\end{enumerate}

\begin{table}[htbp]
\centering
\small
\caption{与本章相关的年份（按时间顺序）}
\label{tab:timeline}
\begin{tabular}{@{}c >{\raggedright\arraybackslash}p{11.2cm}@{}}
\toprule
年份 & 事件 \\
\midrule
1986 & Brown--Henneaux：AdS$_3$ 的渐近对称代数给 $c=3l/2G$（卷页未核）\\
1988 & 本文：CS 理论、共形块态空间对应；“三维引力 $=$ CS”只作为动机出现 \src{396--397}\\
1989 & Moore--Seiberg：Verlinde 猜想的物理证明（见第 \ref{sec:verlinde} 讲）\\
1991 & Reshetikhin--Turaev：量子群与链环不变量的组合构造 \\
1994 & Faltings：Verlinde 公式的代数几何证明 \\
1997 & Maldacena：AdS/CFT 对应（hep-th/9711200，卷页未核）\\
1998 & GKPW：全息字典的精确形式（卷页未核）\\
\bottomrule
\end{tabular}
\end{table}

\subsection{补出的推导}

\begin{derivation}[D-10.1：$Z(X\times S^1)=\dim\cH_X$ 从哪里来]
把 $X\times S^1$ 看成 $X$ 上的“时间演化一个周期”。拓扑场论没有局部自由度，哈密顿量在物理态上为零：
\[
H_X=0 .
\]
理由是：若 $H_X$ 有非零的能隙，配分函数 $Z(X\times S^1)=\Tr_{\cH_X}(e^{-\beta H_X})$ 会依赖 $S^1$ 的长度 $\beta$，这与“不依赖度规”矛盾；拓扑不变性要求 $S^1$ 的长度可以任意缩放而不改变配分函数，故所有物理态的能量必须为零。
于是演化算符是恒等：
\[
U(\beta)=e^{-\beta H_X}=\mathbf 1 ,
\]
从而
\[
Z(X\times S^1)=\Tr_{\cH_X}(U)=\Tr_{\cH_X}(\mathbf 1)=\dim\cH_X .
\]
取 $X=S^2$（无标记点），$\cH_{S^2}=\CC$，得 $Z(S^2\times S^1)=1$，这就是 (4.31)。取 $X=S^2$ 带三个电荷，$\cH_{S^2;R_a,R_b,R_c}$ 的维数就是 $N_{abc}$，这就是 (4.36)。
反过来，这个推导也说明：$X\times S^1$ 上的配分函数\emph{只能}给出维数，不能给出“边界 CFT 的配分函数”。后者会依赖边界上的度量或模参数，而这里没有边界、没有度规。\ours
\end{derivation}

\begin{derivation}[D-10.2：为什么这是“等同”而不是“对偶”]
用模函子的语言把两边写出来。三维 CS 理论给每个定向闭曲面 $\Sigma$（带标记点与表示）一个有限维向量空间 $\cH_\Sigma$，给每个三维配分函数一个数，并对缝合是乘性的。二维 RCFT 给每个 $\Sigma$ 一个共形块空间 $V_\Sigma$，对缝合也是乘性的。原文 p.370 的 “are” 就是说
\[
\cH_\Sigma=V_\Sigma
\]
作为向量空间（连同映射类群作用、Hermite 内积与射影平坦联络）等同。这不是两个理论之间的映射，而是同一个模函子的两种构造。
一个直接推论：三维侧没有“体局部算子”。若有，插入体算子会给出一个新的三维不变量，它不能写成 $\dim\cH_X$ 的形式；本文的构造里没有这样的算子。全息对偶恰恰要求体侧有丰富的局部算子（体场），并与边界算子一一对应。两者结构不同：这里是“一个理论的两种表述”，那里是“两个理论的对偶”。
需要强调：这个“等同”只在态空间层面成立；本文没有断言 $S^3$ 上的配分函数等于某个二维配分函数。$Z(S^3)=S_{0,0}$ 仍是由 $\cH_{T^2}$ 与手术拼出来的数，不是边界 CFT 在某个模参数下的配分函数。\ours
\end{derivation}

\begin{derivation}[D-10.3：$c=3l/2G$ 的量纲检验，以及为什么不能与本文的 $c$ 混同]
三维 Einstein--Hilbert 作用量（带负宇宙常数 $\Lambda=-1/l^2$）是
\[
S_{\mathrm{EH}}=\frac{1}{16\pi G}\int\dd^3x\,\sqrt g\,(R-2\Lambda).
\]
量纲：$[R]=L^{-2}$，$[\dd^3x]=L^3$，$[\sqrt g]=1$，故 $[S_{\mathrm{EH}}]=[G]^{-1}L$。作用量无量纲，所以
\[
[G]=L .
\]
于是 $l/G$ 无量纲（$[l]=L$），中心荷必须形如
\[
c=\alpha\,\frac{l}{G},
\]
其中 $\alpha$ 是无量纲数。Brown--Henneaux 通过求解 AdS$_3$ 的渐近对称代数（两个 Virasoro 代数）定出 $\alpha=3/2$，即 $c=3l/2G$。本讲\emph{不}推导 $3/2$：那需要渐近展开与 Brown--Henneaux 的边界条件计算，超出原文范围。
与本文的中心荷对比：仿射流代数的中心荷是
\[
c=\frac{k\dim G}{k+h^\vee},
\]
$k\in\ZZ$ 离散，$k\to\infty$ 时 $c\to\dim G$ \src{361}。它来自能级 $k$ 的表示论，与 $l/G$ 无关。两个 $c$ 都是无量纲数，但一个是离散的、由仿射代数决定，一个是连续的、由渐近对称代数决定。给定 $k$ 后总能选一个 $l/G$ 使两个数相等，但这只是数值巧合，不是理论等同。\ours
\end{derivation}

\begin{derivation}[D-10.4：“三维引力 $=$ CS”为什么只是经典改写]
三维引力的运动方程是 $R_{\mu\nu}-\frac12Rg_{\mu\nu}+\Lambda g_{\mu\nu}=0$。在 3 维，Weyl 张量恒为零，故真空解局部平坦（$\Lambda=0$）或常曲率（$\Lambda<0$）。把自旋联络与标架组合成 $\mathfrak{iso}(2,1)$（$\Lambda=0$）或 $\mathfrak{sl}(2,\mathbb{R})\oplus\mathfrak{sl}(2,\mathbb{R})$（$\Lambda<0$）值的联络 $A$，Einstein 方程等价于
\[
F(A)=0 ,
\]
即 $A$ 是平坦联络。这正是 CS 理论的驻点方程。所以在\emph{经典}层面，三维引力与取非紧群 $ISO(2,1)$ 或 $SL(2,\mathbb{R})\times SL(2,\mathbb{R})$ 的 CS 理论有相同的解空间。
但“相同解空间”不等于“相同量子理论”，原因有三条。
\begin{enumerate}[label=(\alph*),leftmargin=*]
  \item \textbf{群非紧。}本文的有限维 $\cH_\Sigma$ 论证（p.366--367）依赖 $G$ 紧。非紧群的平坦联络模空间有连续参数，辛商不紧，量子化出现连续谱，不能照搬“有限维”结论。
  \item \textbf{边界条件。}AdS$_3$ 的物理态要满足渐近边界条件；不同的边界条件给出不同的理论（甚至不同的 $c$）。本文在闭流形上工作，没有这一层数据。
  \item \textbf{动力学。}三维引力有边界自由度（渐近 Virasoro、BTZ 黑洞），CS 理论（紧群、闭流形）没有。经典解空间的等同不保证量子配分函数等同；这需要额外的边界项、边界条件与全息字典。
\end{enumerate}
因此，“三维引力 $=$ CS”是一个\emph{经典改写}：它把引力方程翻译成平坦联络条件，是本文与后续全息研究的接触点（原文 p.396--397 与文献 [38]），但不是“本文的 CS 理论 = 量子引力”的证明。\ours
\end{derivation}

\begin{derivation}[D-10.5：为什么本文没有“边界配分函数”可对应]
设 $M$ 是闭三维流形，沿曲面 $\Sigma$ 剪成 $M_1,M_2$。TQFT 公理给出
\[
Z(M)=\langle \psi_1,\psi_2\rangle_{\cH_\Sigma},
\]
其中 $\psi_i\in\cH_\Sigma$ 由 $M_i$ 的“半边”给出。$\cH_\Sigma$ 只依赖 $\Sigma$ 的拓扑（亏格、标记点、表示），不依赖 $\Sigma$ 上的任何共形结构或度量。全息对偶要求的却是另一种等式：
\[
Z_{\text{体}}[X;\text{边界条件}]=Z_{\text{边界 CFT}}[\partial X;\text{边界耦合}],
\]
右端依赖边界上的共形结构（模参数、耦合常数），左端依赖体度规的渐近行为。本文的 $\cH_\Sigma$ 没有模参数可依赖，所以不存在“边界 CFT 配分函数”这一边的对象；能对的只有“维数”，也就是把边界 CFT 的 Hilbert 空间直接取成 $\cH_\Sigma$ 本身。这正是态空间等同与动力学等价的分界线。\ours
\end{derivation}

\begin{derivation}[D-10.6：等同的对象只能是共形块，不是整个 CFT Hilbert 空间]
取 $X=T^2$。三维 CS 的态空间维数是 $t$（$SU(2)$ 时 $k+1$），有限维（原文 \S4.3 的 (4.24)--(4.26)；\src{383--384}）。而任何一个二维 CFT 在环面上的 Hilbert 空间都是无穷维的：固定中心荷与模参数后，每个共形族仍有整族的 Virasoro 后代 $L_{-n_1}\cdots L_{-n_s}|h\rangle$，$s$ 任意，故态数无限。两边的维数一个有限、一个无限，所以
\[
\cH^{\text{3d}}_{T^2}\;\cong\;\text{（固定标记点与表示下的共形块空间）}
\]
只能是“共形块子空间”层面的等同，不可能是“整个二维 CFT Hilbert 空间”的等同。这也从态空间一侧解释了为什么本文没有全息字典：字典要把体态与边界态一一对应，而体侧在 $T^2$ 上只有 $t$ 个态，边界侧有无穷多个；比例 $t/\infty$ 说明对应只能落在共形块上。反过来，共形块空间正是共形场论里“由 Ward 恒等式与表示论决定、不依赖动力学细节”的那部分，所以这个等同虽窄，却是精确的。\ours
\end{derivation}

\begin{derivation}[D-10.7：三维作用量不含度规，所以没有“对偶的边界应力张量”]
在 AdS/CFT 里，边界应力张量由对边界度规变分定义：
\[
T_{ij}=\frac{2}{\sqrt h}\frac{\delta S_{\text{ren}}}{\delta h^{ij}},
\]
它的两点函数给出中心荷。三维 CS 的作用量 (1.3) 只用到流形上的定向与丛，不含任何度规；把边界度规 $h_{ij}$ 当参数变分，
\[
\frac{\delta S_{\text{CS}}}{\delta h^{ij}}=0
\]
是恒等式。于是三维侧没有任何算符与 $T_{ij}$ 对偶，也就没有“体场 $\leftrightarrow$ 边界应力张量”这一条字典。
二维侧当然有 Sugawara 应力张量
\[
T(z)=\frac{1}{2(k+h^\vee)}\sum_a {:}J^aJ^a{:}(z),
\]
它给出 Virasoro 代数与 $c=k\dim G/(k+h^\vee)$；但它的存在是二维共形场论自己的性质，不需要三维侧的度规来“源”它。两个 $c$ 的数值可能在某些 $k,l,G$ 下恰好相等（见练习 6），但产生机制不同：一个来自二维流的 Sugawara 构造，一个来自三维度规的渐近对称分析。这正是“态空间等同”与“动力学等价”的分水岭。\ours
\end{derivation}

\subsection{与其他来源的对照}

\begin{center}
\begin{minipage}{\linewidth}
\centering
\captionof{table}{两种“三维--二维关系”的逐项对照}
\label{tab:duality}
\small
\small
\setlength{\tabcolsep}{3pt}
\begin{tabular}{@{}>{\raggedright\arraybackslash}p{2.4cm} >{\raggedright\arraybackslash}p{6.6cm} >{\raggedright\arraybackslash}p{6.2cm}@{}}
\toprule
项目 & CS/WZW（本文，1988） & AdS$_3$/CFT$_2$（1997 后）\\
\midrule
体理论 & 三维 Chern--Simons，作用量 (1.3) 不含度规 \src{354} & 三维量子引力，含度规与负宇宙常数 \\
体自由度 & 无；$\cH_\Sigma$ 有限维 \src{366--367} & 有：引力子、BTZ 黑洞、黑洞熵 \\
边界条件 & 由 $\Sigma$ 的拓扑给定（亏格、标记点），无渐近条件 & 固定渐近 AdS$_3$ 度规，两个 Virasoro 渐近对称 \\
关系类型 & 态空间等同：$\cH_\Sigma=V_\Sigma$（p.370 用 “are”）\src{370} & 动力学等价：$Z_{\text{体}}=Z_{\text{边界}}$，靠全息字典 \\
三维接触二维的方式 & 仅 $Z(X\times S^1)=\dim\cH_X$ \src{387--388} & 一般 $X$：体配分函数 = 边界配分函数 \\
中心荷 & $c=k\dim G/(k+h^\vee)$ \src{390}；离散，$k\in\ZZ$ \src{354} & $c=3l/2G$（Brown--Henneaux）；连续 \\
可证性 & 可逐条检验（几何量子化、Narasimhan--Seshadri、BWB） & 无定理；只有弱耦合、对称性、可积性检验 \\
全息字典 & 无 & 有（体场 $\leftrightarrow$ 边界算子，质量 $\leftrightarrow$ 共形维数）\\
\bottomrule
\end{tabular}
\end{minipage}
\end{center}

\begin{center}
\begin{minipage}{\linewidth}
\centering
\captionof{table}{各来源对“3d/2d 关系”的表述与核对程度}
\label{tab:3d2dsources}
\small
\small
\setlength{\tabcolsep}{3pt}
\begin{tabular}{@{}>{\raggedright\arraybackslash}p{2.7cm} >{\raggedright\arraybackslash}p{6.6cm} >{\raggedright\arraybackslash}p{2.9cm} >{\raggedright\arraybackslash}p{2.8cm}@{}}
\toprule
来源 & 对 3d/2d 关系的表述 & 是否涉及全息 & 本讲核对程度 \\
\midrule
本文 & “共形块空间\emph{是}量子化三维理论得到的 Hilbert 空间”（p.370 用 “are”） & 否，全文零命中 & 逐页核对 \\
河野 \S2.3--2.4 & 作为 CS/WZW 对应的背景（书 p.58、p.63） & 未讨论 & 仅据研讨记录，未逐字核 \\
上野 \S5 & 讲融合规则与 Verlinde 公式（p.101--107） & 未讨论 & 未逐字核 \\
Moore \S4 & 讲二维 RCFT 的融合与模性质 & 未讨论 & 未核 \\
Schottenloher 第 11 章 & 数学侧 Verlinde 与共形块（p.171--178），不涉及三维引力 & 未讨论 & p.171--178 已核 \\
李思 & 未核 & —— & 未核 \\
\bottomrule
\end{tabular}
\end{minipage}
\end{center}

表 \ref{tab:duality} 的第三列全部属于 1997 年之后的框架；把它与第二列并排放，是为了让“结构差异”一眼可见，不是声称本文预言了第三列。尤其“中心荷”一行：$k\dim G/(k+h^\vee)$ 与 $3l/2G$ 都是无量纲数，但前者是仿射代数的能级公式，后者是三维渐近对称代数的中心荷，来源不同、参数类型不同（离散 vs 连续），不能互相代入。

表 \ref{tab:3d2dsources} 里出现若干“未核”，这是刻意的：本讲的核对以原文（逐页）与研讨记录（S5、S6 报告）为限；二手文献只报页码与结论，不冒充逐字引用。读者若要引用河野、上野、Moore 的具体表述，请以原书为准。反过来，凡是标“逐页核对”的，都是把英文原句与页码一起给出的，可以直接复核。

还有一条差别值得单独说：可证性。CS/WZW 的态空间等同有定理级的支撑——Segal 的模函子定义、TUY 的共形块丛构造，加上 Verlinde 公式的两条独立数学证明（第 9 讲表 \ref{tab:routes}）；三维侧每一步都可以写成“某空间 = 某空间”并逐项核对。AdS$_3$/CFT$_2$ 至今没有同等级的定理：它依赖大 $c$、弱耦合与对称性假设，检验方式是匹配配分函数、谱与关联函数，而不是证明两个模函子相同。因此“CS/WZW 是严格结果，全息是框架”不是价值判断，而是两类陈述的逻辑地位不同。这也是本讲把“态空间等同”与“动力学等价”分开的最后一条理由。\ours

最后把三句话再说一遍，作为本章的结论。第一，本文的三维--二维关系是\emph{态空间等同}：同一个模函子，两种构造，原文用 “are”。第二，AdS$_3$/CFT$_2$ 是\emph{动力学等价}：两个理论、边界条件、全息字典，$c=3l/2G$ 属于这一边。第三，“三维引力 $=$ CS”是两者之间唯一的\emph{经典}接触点：非紧群、只等同解空间，既不是本文的定理，也不是全息的雏形。

\subsection{陷阱与约定}

\begin{keypoint}
\textbf{三种关系必须分开叫：}(i) CS/WZW 是\emph{态空间等同}，原文用 “are”，只在 $X\times S^1$ 上以 $Z=\dim\cH_X$ 接触二维数据；(ii) AdS$_3$/CFT$_2$ 是\emph{动力学等价}，需要边界条件、大 $c$ 与全息字典；(iii) “三维引力 $=$ CS”是\emph{经典改写}，用非紧群、只等同解空间。把三者中的任意两个混为一谈，都会得出“本文证明了全息”或“本文的 CS 就是量子引力”这类错误结论。
\end{keypoint}

\begin{itemize}[leftmargin=1.6em]
  \item \textbf{“are” 不是 “correspond”。}p.370 的动词选择是有意的：同一个向量空间，不是两个空间之间的同构映射（虽然可以等价地这么表述）。把它读成“体配分函数 = 边界配分函数”是过度解读。
  \item \textbf{不要用本文“证明”全息。}本文既没有度规，也没有边界条件与字典；$c=3l/2G$ 不在本文中出现。用 1988 年的文章支持 1997 年的猜想，是时代错置。
  \item \textbf{$3d$ 引力的规范群。}$\Lambda=0$ 用 $ISO(2,1)$ 或其覆盖群；$\Lambda<0$ 用 $SL(2,\mathbb{R})\times SL(2,\mathbb{R})$。两者都\emph{非紧}，与本文的紧单连通 $G$ 不同。原文只在 p.396--397 一句带过，细节在 [38]。
  \item \textbf{两个 $c$ 不能混。}本文 $c=k\dim G/(k+h^\vee)$ 是仿射中心荷，离散；$c=3l/2G$ 是渐近 Virasoro 中心荷，连续。数值上可调参相等，理论上不等同。
  \item \textbf{有限维依赖紧性。}p.366--367 的有限维论证对紧群成立；非紧群的平坦联络模空间不紧，不能照搬。不要把“有限维 $\cH_\Sigma$”当成所有三维引力理论的普遍性质。
  \item \textbf{时间线。}1988：本文只有态空间对应与一句“三维引力=CS”的动机；1997：Maldacena 提出 AdS/CFT；1998：GKPW 给出字典。按这条时间线读，本文是“前史”，不是“雏形”。
\end{itemize}

\subsection{练习}

\begin{enumerate}[leftmargin=*]
  \item （复算）用 (4.30) 验证 $Z(S^2\times S^1)=1$，并用 (4.34) 验证单电荷时为 $\delta_{a0}$。说明为什么这两个等式不能读成“边界 CFT 的配分函数”。
  \item （推导）从 $H_X=0$ 推出 $Z(X\times S^1)=\dim\cH_X$；若 $H_X$ 有能隙 $\Delta>0$，$Z$ 会变成什么？为什么拓扑不变性排除了这种情形？
  \item （推导）在三维中验证 $[G]=L$，并由此说明 $c=3l/2G$ 无量纲。为什么量纲分析只能定出 $c=\alpha\,l/G$，定不出 $\alpha=3/2$？
  \item （判断）“CS/WZW 是 AdS$_3$/CFT$_2$ 全息对偶在 1988 年的雏形。”这句话对吗？请从度规、边界条件、全息字典三个方面各给一条反驳。
  \item （判断）“本文证明了三维引力可以写成 Chern--Simons 理论，所以本文的 CS 理论就是量子引力。”错在哪里？请指出非紧群、经典层面、边界条件三个关键词。
  \item （判断）取 $SU(2)$、$k=1$，则 $c=k\dim G/(k+h^\vee)=1$。能否选 $l/G=2/3$ 使 $c=3l/2G=1$，从而“证明”两个理论相同？请说明数值相等为什么不构成理论等同。
  \item （复算）用 (4.39) 算 $k=1$ 时的 $Z(S^3)=S_{0,0}$，并说明它是 $\cH_{T^2}$ 上的一个数，而不是某个边界 CFT 的配分函数。
  \item （判断）$SU(2)_k$ 的 $\dim\cH_{T^2}=k+1$ 是有限的，而二维 CFT 在环面上的 Hilbert 空间是无穷维的。由此说明：本文的“等同”为什么只能发生在共形块层面？这个论证能否反过来当作“本文不是全息”的证据？
\end{enumerate}

\begin{boundary}
\begin{itemize}[leftmargin=1.4em]
  \item \textbf{原文证明：}作用量 (1.3) 不含度规 \src{353--354}；$\cH_\Sigma$ 有限维 \src{366--367}；共形块空间“是”量子化三维理论所得的 Hilbert 空间 \src{370}；$Z(X\times S^1)=\dim\cH_X$ \src{387--388}；$Z(S^3)=S_{0,0}$ \src{389}。原文对“三维引力=CS”只有一句动机 \src{396--397} 与文献 [38]。
  \item \textbf{我们补的：}D-10.1 的 $H=0$ 推导；D-10.2 的模函子表述；D-10.3 的量纲检验与两个 $c$ 的区分；D-10.4 的经典改写论证与三条障碍；表 \ref{tab:duality} 的八项对照；表 \ref{tab:3d2dsources} 的来源核对。
  \item \textbf{未证／未核：}Brown--Henneaux 的 $c=3l/2G$（1986；题名已核，卷页未核）；Achúcarro--Townsend 的非紧 CS 表述（1986，卷页未核）；Maldacena 1997（hep-th/9711200，卷页未核）、GKPW 1998、BTZ 1992 的卷页均未核；非紧群的量子化与连续谱只给出定性论证，未作严格分析（Witten 2007 的讨论未核）；河野 p.58、p.63、上野 p.101--107、Moore \S4、李思的对应表述未逐字核对。
  \item \textbf{练习答案：}本讲只给提示，未写入附录 C（本讲的写权限不含 \texttt{appC}）。
\end{itemize}
\end{boundary}
